% Options for packages loaded elsewhere
\PassOptionsToPackage{unicode}{hyperref}
\PassOptionsToPackage{hyphens}{url}
%
\documentclass[
]{article}
\usepackage{amsmath,amssymb}
\usepackage{iftex}
\ifPDFTeX
  \usepackage[T1]{fontenc}
  \usepackage[utf8]{inputenc}
  \usepackage{textcomp} % provide euro and other symbols
\else % if luatex or xetex
  \usepackage{unicode-math} % this also loads fontspec
  \defaultfontfeatures{Scale=MatchLowercase}
  \defaultfontfeatures[\rmfamily]{Ligatures=TeX,Scale=1}
\fi
\usepackage{lmodern}
\ifPDFTeX\else
  % xetex/luatex font selection
\fi
% Use upquote if available, for straight quotes in verbatim environments
\IfFileExists{upquote.sty}{\usepackage{upquote}}{}
\IfFileExists{microtype.sty}{% use microtype if available
  \usepackage[]{microtype}
  \UseMicrotypeSet[protrusion]{basicmath} % disable protrusion for tt fonts
}{}
\makeatletter
\@ifundefined{KOMAClassName}{% if non-KOMA class
  \IfFileExists{parskip.sty}{%
    \usepackage{parskip}
  }{% else
    \setlength{\parindent}{0pt}
    \setlength{\parskip}{6pt plus 2pt minus 1pt}}
}{% if KOMA class
  \KOMAoptions{parskip=half}}
\makeatother
\usepackage{xcolor}
\usepackage{longtable,booktabs,array}
\usepackage{calc} % for calculating minipage widths
% Correct order of tables after \paragraph or \subparagraph
\usepackage{etoolbox}
\makeatletter
\patchcmd\longtable{\par}{\if@noskipsec\mbox{}\fi\par}{}{}
\makeatother
% Allow footnotes in longtable head/foot
\IfFileExists{footnotehyper.sty}{\usepackage{footnotehyper}}{\usepackage{footnote}}
\makesavenoteenv{longtable}
\usepackage{graphicx}
\makeatletter
\def\maxwidth{\ifdim\Gin@nat@width>\linewidth\linewidth\else\Gin@nat@width\fi}
\def\maxheight{\ifdim\Gin@nat@height>\textheight\textheight\else\Gin@nat@height\fi}
\makeatother
% Scale images if necessary, so that they will not overflow the page
% margins by default, and it is still possible to overwrite the defaults
% using explicit options in \includegraphics[width, height, ...]{}
\setkeys{Gin}{width=\maxwidth,height=\maxheight,keepaspectratio}
% Set default figure placement to htbp
\makeatletter
\def\fps@figure{htbp}
\makeatother
\setlength{\emergencystretch}{3em} % prevent overfull lines
\providecommand{\tightlist}{%
  \setlength{\itemsep}{0pt}\setlength{\parskip}{0pt}}
\setcounter{secnumdepth}{-\maxdimen} % remove section numbering
\ifLuaTeX
  \usepackage{selnolig}  % disable illegal ligatures
\fi
\IfFileExists{bookmark.sty}{\usepackage{bookmark}}{\usepackage{hyperref}}
\IfFileExists{xurl.sty}{\usepackage{xurl}}{} % add URL line breaks if available
\urlstyle{same}
\hypersetup{
  hidelinks,
  pdfcreator={LaTeX via pandoc}}

\author{}
\date{}

\begin{document}


\includegraphics[width=5.54583in,height=0.92431in]{media/image1.png}


\includegraphics[width=3.68194in,height=0.73681in]{media/image2.png}

\textbf{ÉCOLE SUPÉRIEURE DE TECHNOLOGIE - FÈS}

\textbf{RAPPORT DE STAGE}

\textbf{DIPLÔME BACHELOR EN TECHNOLOGIE}

\textbf{FILIÈRE : Génie LOGICIELLE (GL)}

\emph{\textbf{Conception et Développement d\textquotesingle une
Plateforme Intelligente de Présélection des Candidats}}

\textbf{Réalisé par :}

\begin{itemize}
\item
  \textbf{SAMRANI OUAHID}
\end{itemize}

\textbf{Soutenu devant le jury :}

\textbf{- YOUNES AZIZ~BACHIRI (Encadrant académique)}

\textbf{- DAHBI MAROUANE (Encadrant entreprise)}

Lieu de stage : ENSA Béni Mellal --- USMS

\begin{quote}
06/04/2026 au 13/06/2026 (Année Universitaire : 2025-2026)
\end{quote}

\hypertarget{remerciements}{%
\section{\texorpdfstring{Remerciements
}{Remerciements }}\label{remerciements}}

Avant d\textquotesingle entamer la présentation de ce travail, je tiens
à exprimer ma profonde gratitude envers toutes les personnes qui ont
contribué à la réussite de ce stage.

Je remercie tout d\textquotesingle abord mon encadrant professionnel au
sein de l\textquotesingle ENSA Béni Mellal, \textbf{DAHBI MAROUANE} ,
pour son encadrement rigoureux, ses orientations pertinentes et sa
disponibilité tout au long de la période de stage.

J\textquotesingle adresse également mes sincères remerciements à mon
encadrant académique à l\textquotesingle EST de Fès, \textbf{YOUNES
AZIZ~BACHIRI}, pour son suivi pédagogique et ses recommandations
précieuses.

Ma reconnaissance va aussi à l\textquotesingle ensemble du personnel du
service de la Scolarité de l\textquotesingle ENSA BM, qui a facilité mon
intégration et m\textquotesingle a fourni les informations nécessaires à
la compréhension du processus de présélection.

Je remercie les membres du jury, \textbf{YOUNES AZIZ~BACHIRI},
d\textquotesingle avoir accepté d\textquotesingle évaluer ce travail.

Enfin, je remercie l\textquotesingle USMS et l\textquotesingle EST de
Fès pour la qualité de la formation dispensée en Génie Logiciel.

\hypertarget{ruxe9sumuxe9}{%
\section{Résumé}\label{ruxe9sumuxe9}}

Le présent rapport rend compte du stage de fin d\textquotesingle études
effectué au sein de l\textquotesingle ENSA Béni Mellal, rattachée à
l\textquotesingle Université Sultan Moulay Slimane. Ce stage
s\textquotesingle inscrit dans le cadre de l\textquotesingle obtention
du Bachelor Universitaire de Technologie en Génie Logiciel de
l\textquotesingle EST de Fès.

Le projet consiste en la conception et le développement
d\textquotesingle une plateforme web intelligente dédiée à la
présélection des candidats aux filières d\textquotesingle ingénieur de
l\textquotesingle ENSA BM. Cette plateforme automatise le traitement des
dossiers de candidature, depuis le dépôt en ligne
jusqu\textquotesingle au classement final, en passant par la
vérification OCR des documents, le scoring académique pondéré par
filière et la gestion des épreuves écrites.

La solution repose sur une architecture Full Stack moderne : Django REST
Framework pour le backend et React 18 avec Tailwind CSS pour le
frontend. Elle intègre un moteur de règles de rejet automatique (7
types), un algorithme de scoring avec fuzzy matching des matières, un
système de notifications en temps réel via WebSocket, et un module
d\textquotesingle import des notes d\textquotesingle épreuve écrite au
format Excel.

\textbf{Mots-clés : Présélection, ENSA, Django, React, Scoring
automatique, OCR, Fuzzy Matching, WebSocket, LMD, Full Stack.}

\hypertarget{abstract}{%
\subsection{Abstract}\label{abstract}}

This report presents the internship carried out at the National School
of Applied Sciences of Béni Mellal (ENSA BM), part of Sultan Moulay
Slimane University. The project involves the design and development of
an intelligent web platform for candidate preselection to ENSA BM
engineering programs.

The solution is built on a modern Full Stack architecture combining
Django REST Framework for the backend and React 18 with Tailwind CSS for
the frontend. It features an automatic rejection rule engine (7 types),
a weighted scoring algorithm with fuzzy matching, real-time WebSocket
notifications, and an Excel-based exam score import module.

\textbf{Keywords: Preselection, ENSA, Django, React, Automatic scoring,
OCR, Fuzzy Matching, WebSocket, LMD, Full Stack.}

\hypertarget{sommaire}{%
\section{Sommaire}\label{sommaire}}

\hypertarget{sommaire-1}{%
\section{Sommaire}\label{sommaire-1}}

\protect\hyperlink{remerciements}{Remerciements 2}

\protect\hyperlink{ruxe9sumuxe9}{Résumé 3}

\protect\hyperlink{abstract}{Abstract 3}

\protect\hyperlink{sommaire}{Sommaire 5}

\protect\hyperlink{__RefHeading___Toc5287_1548811319}{Liste des figures
10}

\protect\hyperlink{__RefHeading___Toc5289_1548811319}{Liste des tableaux
11}

\protect\hyperlink{liste-des-abruxe9viations}{Liste des abréviations 12}

\protect\hyperlink{introduction-guxe9nuxe9rale}{Introduction générale
14}

\protect\hyperlink{chapitre-1-contexte-du-projet}{Chapitre 1 : Contexte
du Projet 15}

\protect\hyperlink{introduction}{1.1 Introduction 15}

\protect\hyperlink{pruxe9sentation-de-lorganisme-daccueil}{1.2
Présentation de l\textquotesingle organisme d\textquotesingle accueil
15}

\protect\hyperlink{lensa-buxe9ni-mellal}{1.2.1 L\textquotesingle ENSA
Béni Mellal 15}

\protect\hyperlink{missions-et-offre-de-formation}{1.2.2 Missions et
offre de formation 16}

\protect\hyperlink{service-de-la-scolarituxe9}{1.2.3 Service de la
Scolarité 16}

\protect\hyperlink{contexte-et-probluxe9matique}{1.3 Contexte et
problématique 17}

\protect\hyperlink{le-processus-actuel-de-pruxe9suxe9lection}{1.3.1 Le
processus actuel de présélection 17}

\protect\hyperlink{probluxe8mes-identifiuxe9s}{1.3.2 Problèmes
identifiés 17}

\protect\hyperlink{le-systuxe8me-lmd-marocain}{1.3.3 Le système LMD
marocain 17}

\protect\hyperlink{objectifs-du-projet}{1.4 Objectifs du projet 18}

\protect\hyperlink{objectif-principal}{1.4.1 Objectif principal 18}

\protect\hyperlink{objectifs-spuxe9cifiques}{1.4.2 Objectifs spécifiques
18}

\protect\hyperlink{puxe9rimuxe8tre-fonctionnel}{1.5 Périmètre
fonctionnel 18}

\protect\hyperlink{muxe9thodologie-de-travail}{1.6 Méthodologie de
travail 19}

\protect\hyperlink{conclusion}{1.7 Conclusion 20}

\protect\hyperlink{chapitre-2-architecture-et-base-de-donnuxe9es}{Chapitre
2 : Architecture et Base de données 21}

\protect\hyperlink{introduction-1}{2.1 Introduction 21}

\protect\hyperlink{architecture-globale}{2.2 Architecture globale 22}

\protect\hyperlink{vue-densemble}{2.2.1 Vue d\textquotesingle ensemble
22}

\protect\hyperlink{architecture-rest-stateless}{2.2.2 Architecture REST
stateless 24}

\protect\hyperlink{stack-technologique}{2.3 Stack technologique 24}

\protect\hyperlink{organisation-modulaire-du-backend}{2.4 Organisation
modulaire du backend 25}

\protect\hyperlink{dictionnaire-des-moduxe8les-django}{2.5 Dictionnaire
des modèles Django 27}

\protect\hyperlink{application-users}{2.5.1 Application Users 29}

\protect\hyperlink{application-candidatures}{2.5.2 Application
Candidatures 30}

\protect\hyperlink{application-administration-et-scoring}{2.5.3
Application Administration et Scoring 31}

\protect\hyperlink{relations-et-cardinalituxe9s}{2.6 Relations et
cardinalités 32}

\protect\hyperlink{machine-duxe9tats-du-dossier}{2.7 Machine
d\textquotesingle états du dossier 33}

\protect\hyperlink{conclusion-1}{2.8 Conclusion 34}

\protect\hyperlink{chapitre-3-logique-muxe9tier-et-algorithmes}{Chapitre
3 : Logique Métier et Algorithmes 36}

\protect\hyperlink{introduction-2}{3.1 Introduction 36}

\protect\hyperlink{moteur-de-ruxe8gles-de-rejet-automatique}{3.2 Moteur
de règles de rejet automatique 36}

\protect\hyperlink{principe-de-fonctionnement}{3.2.1 Principe de
fonctionnement 36}

\protect\hyperlink{types-de-ruxe8gles-impluxe9mentuxe9s}{3.2.2 Types de
règles implémentés 36}

\protect\hyperlink{journalisation-des-rejets}{3.2.3 Journalisation des
rejets 37}

\protect\hyperlink{algorithme-de-scoring-ponduxe9ruxe9-par-filiuxe8re}{3.3
Algorithme de scoring pondéré par filière 37}

\protect\hyperlink{principe-guxe9nuxe9ral}{3.3.1 Principe général 37}

\protect\hyperlink{coefficients-des-quatre-filiuxe8res}{3.3.2
Coefficients des quatre filières 38}

\protect\hyperlink{formule-de-calcul-du-score-dossier}{3.3.3 Formule de
calcul du score dossier 38}

\protect\hyperlink{bonus-de-mention-et-coefficient-dipluxf4me}{3.3.4
Bonus de mention et coefficient diplôme 39}

\protect\hyperlink{formule-du-score-final}{3.3.5 Formule du score final
40}

\protect\hyperlink{pipeline-ocr-et-fuzzy-matching}{3.4 Pipeline OCR et
Fuzzy Matching 40}

\protect\hyperlink{extraction-ocr}{3.4.1 Extraction OCR 40}

\protect\hyperlink{nettoyage-du-bruit-ocr}{3.4.2 Nettoyage du bruit OCR
41}

\protect\hyperlink{fuzzy-matching-des-matiuxe8res}{3.4.3 Fuzzy Matching
des matières 42}

\protect\hyperlink{scoring-avec-fallback}{3.4.4 Scoring avec fallback
42}

\protect\hyperlink{calcul-automatique-des-mentions}{3.5 Calcul
automatique des mentions 43}

\protect\hyperlink{conclusion-2}{3.6 Conclusion 43}

\protect\hyperlink{chapitre-4-impluxe9mentation-asynchronisme-et-tests}{Chapitre
4 : Implémentation, Asynchronisme et Tests 44}

\protect\hyperlink{introduction-3}{4.1 Introduction 44}

\protect\hyperlink{interfaces-duxe9veloppuxe9es}{4.2 Interfaces
développées 44}

\protect\hyperlink{gestion-de-lasynchronisme}{4.3 Gestion de
l\textquotesingle asynchronisme 49}

\protect\hyperlink{probluxe9matique}{4.3.1 Problématique 49}

\protect\hyperlink{solution-threading.thread}{4.3.2 Solution :
threading.Thread 49}

\protect\hyperlink{gestion-des-connexions-base-de-donnuxe9es}{4.3.3
Gestion des connexions base de données 49}

\protect\hyperlink{systuxe8me-de-notifications-en-temps-ruxe9el}{4.4
Système de notifications en temps réel 49}

\protect\hyperlink{architecture-websocket}{4.4.1 Architecture WebSocket
49}

\protect\hyperlink{envoi-en-masse-avec-progression}{4.4.2 Envoi en masse
avec progression 50}

\protect\hyperlink{systuxe8me-de-logging}{4.5 Système de logging 50}

\protect\hyperlink{suxe9curituxe9-de-lapplication}{4.6 Sécurité de
l\textquotesingle application 50}

\protect\hyperlink{stratuxe9gie-de-tests}{4.7 Stratégie de tests 51}

\protect\hyperlink{niveaux-de-tests}{4.7.1 Niveaux de tests 51}

\protect\hyperlink{tests-du-moteur-de-ruxe8gles}{4.7.2 Tests du moteur
de règles 51}

\protect\hyperlink{tests-de-lalgorithme-de-scoring}{4.7.3 Tests de
l\textquotesingle algorithme de scoring 52}

\protect\hyperlink{ruxe9capitulatif-des-tests}{4.7.4 Récapitulatif des
tests 52}

\protect\hyperlink{couverture-de-code}{4.7.5 Couverture de code 53}

\protect\hyperlink{difficultuxe9s-rencontruxe9es-et-solutions}{4.8
Difficultés rencontrées et solutions 53}

\protect\hyperlink{conclusion-3}{4.9 Conclusion 54}

\protect\hyperlink{conclusion-guxe9nuxe9rale}{Conclusion générale 55}

\protect\hyperlink{perspectives-damuxe9lioration}{Perspectives
d\textquotesingle amélioration 55}

\protect\hyperlink{bibliographie}{Bibliographie 57}

Table of Figures

\protect\hyperlink{Figureux210ux7csequence}{Figure 1: Diagramme de cas
d\textquotesingle utilisation global de la plateforme 21}

\protect\hyperlink{Figureux211ux7csequence}{Figure 2: Architecture
générale de la plateforme de présélection 22}

\protect\hyperlink{Figureux212ux7csequence}{Figure 3: Organisation
modulaire du backend Django 25}

\protect\hyperlink{Figureux213ux7csequence}{Figure 4: Diagramme de
classes de la plateforme 27}

\protect\hyperlink{Figureux214ux7csequence}{Figure 5: Diagramme de
séquence : Soumission d\textquotesingle une candidature 30}

\protect\hyperlink{Figureux215ux7csequence}{Figure 6: Machine
d\textquotesingle états du dossier de candidature 33}

\protect\hyperlink{Figureux216ux7csequence}{Figure 7: Fonctionnement du
moteur de rejet automatique 36}

\protect\hyperlink{Figureux217ux7csequence}{Figure 8: Processus de
calcul du score pondéré 38}

\protect\hyperlink{Figureux218ux7csequence}{Figure 9: Pipeline OCR et
traitement des relevés de notes 40}

\protect\hyperlink{Figureux219ux7csequence}{Figure 10: Interface
d\textquotesingle accueil de la plateforme 44}

\protect\hyperlink{Figureux2110ux7csequence}{Figure 11: Formulaire
candidature 45}

\protect\hyperlink{Figureux2111ux7csequence}{Figure 12: Interface de
dépôt des documents justificatifs 45}

\protect\hyperlink{Figureux2112ux7csequence}{Figure 13: Tableau de bord
du responsable de présélection 46}

\protect\hyperlink{Figureux2113ux7csequence}{Figure 14: Consultation
détaillée d\textquotesingle un dossier candidat 46}

\protect\hyperlink{Figureux2114ux7csequence}{Figure 15: Configuration
des filières et règles de scoring 47}

\protect\hyperlink{Figureux2115ux7csequence}{Figure 16: Tableau de bord
analytique du responsable de présélection 47}

\hypertarget{section}{%
\section{}\label{section}}

Index of Tables

\protect\hyperlink{Tableux210ux7csequence}{Table 1: Fiche signalétique
de l\textquotesingle ENSA BM 15}

\protect\hyperlink{Tableux211ux7csequence}{Table 2: Filières
d\textquotesingle ingénieur de l\textquotesingle ENSA BM 15}

\protect\hyperlink{Tableux212ux7csequence}{Table 3: Périmètre
fonctionnel du projet 18}

\protect\hyperlink{Tableux213ux7csequence}{Table 4: Planning
prévisionnel du projet 19}

\protect\hyperlink{Tableux214ux7csequence}{Table 5: Stack technologique
du projet 23}

\protect\hyperlink{Tableux215ux7csequence}{Table 6: Applications Django
du backend 24}

\protect\hyperlink{Tableux216ux7csequence}{Table 7: Services métier du
module candidatures 25}

\protect\hyperlink{Tableux217ux7csequence}{Table 8: Modèles
d\textquotesingle administration et scoring 31}

\protect\hyperlink{Tableux218ux7csequence}{Table 9: Relations et
cardinalités du modèle de données 32}

\protect\hyperlink{Tableux219ux7csequence}{Table 10: Tableau 2.6 :
Machine d\textquotesingle états du dossier de candidature 33}

\protect\hyperlink{Tableux2110ux7csequence}{Table 11: Les 7 types de
règles de rejet automatique 36}

\protect\hyperlink{Tableux2111ux7csequence}{Table 12: Matrice des
pondérations par filière 37}

\protect\hyperlink{Tableux2112ux7csequence}{Table 13: Corrections
automatiques du bruit OCR 41}

\protect\hyperlink{Tableux2113ux7csequence}{Table 14: Interfaces
développées par rôle 44}

\protect\hyperlink{Tableux2114ux7csequence}{Table 15: Configuration des
loggers 50}

\protect\hyperlink{Tableux2115ux7csequence}{Table 16: Mesures de
sécurité implémentées 51}

\protect\hyperlink{Tableux2116ux7csequence}{Table 17: Difficultés
rencontrées et solutions 53}

\hypertarget{liste-des-abruxe9viations}{%
\section{Liste des abréviations}\label{liste-des-abruxe9viations}}

\begin{longtable}[]{@{}
  >{\raggedright\arraybackslash}p{(\columnwidth - 2\tabcolsep) * \real{0.5001}}
  >{\raggedright\arraybackslash}p{(\columnwidth - 2\tabcolsep) * \real{0.4999}}@{}}
\toprule\noalign{}
\begin{minipage}[b]{\linewidth}\raggedright
\textbf{Abréviation}
\end{minipage} & \begin{minipage}[b]{\linewidth}\raggedright
\textbf{Signification}
\end{minipage} \\
\midrule\noalign{}
\endhead
\bottomrule\noalign{}
\endlastfoot
ENSA & École Nationale des Sciences Appliquées \\
USMS & Université Sultan Moulay Slimane \\
EST & École Supérieure de Technologie \\
USMBA & Université Sidi Mohamed Ben Abdellah \\
LMD & Licence --- Master --- Doctorat \\
API & Application Programming Interface \\
REST & Representational State Transfer \\
DRF & Django REST Framework \\
JWT & JSON Web Token \\
RBAC & Role-Based Access Control \\
OCR & Optical Character Recognition \\
SMTP & Simple Mail Transfer Protocol \\
CSRF & Cross-Site Request Forgery \\
XSS & Cross-Site Scripting \\
SPA & Single Page Application \\
ORM & Object-Relational Mapping \\
TDI & Télécommunications et Développement Informatique \\
IACS & Informatique Appliquée et Cybersécurité \\
IAA & Ingénierie Agro-Alimentaire \\
G2ER & Génie de l\textquotesingle Énergie et Énergies Renouvelables \\
CIN & Carte d\textquotesingle Identité Nationale \\
CNE & Code National de l\textquotesingle Étudiant \\
KPI & Key Performance Indicator \\
\end{longtable}

\hypertarget{introduction-guxe9nuxe9rale}{%
\section{Introduction générale}\label{introduction-guxe9nuxe9rale}}

Dans un contexte de transformation numérique des établissements
d\textquotesingle enseignement supérieur au Maroc, la gestion des
processus d\textquotesingle admission et de présélection constitue un
enjeu majeur pour les écoles d\textquotesingle ingénieurs.
L\textquotesingle ENSA Béni Mellal accueille chaque année un nombre
croissant de candidatures pour ses filières d\textquotesingle ingénieur,
tant au niveau Bac+2 qu\textquotesingle au niveau Bac+3.

Le processus actuel de présélection, principalement manuel et basé sur
des fichiers Excel, présente plusieurs limites : risque
d\textquotesingle erreurs humaines, lenteur du traitement, absence de
traçabilité et difficulté de gestion des volumes importants de dossiers.
Le service de la Scolarité a exprimé le besoin d\textquotesingle une
solution informatique capable d\textquotesingle automatiser et de
fiabiliser ce processus.

C\textquotesingle est dans ce contexte que s\textquotesingle inscrit
notre projet de stage : la conception et le développement
d\textquotesingle une plateforme web intelligente de présélection des
candidats. Ce rapport est structuré en quatre chapitres :

\begin{itemize}
\item
  Chapitre 1 --- Contexte du Projet : présentation de
  l\textquotesingle organisme d\textquotesingle accueil, problématique
  et objectifs.
\item
  Chapitre 2 --- Architecture et Base de données : architecture
  Full-Stack, stack technologique et dictionnaire des modèles Django.
\item
  Chapitre 3 --- Logique Métier et Algorithmes : moteur de règles,
  scoring pondéré, pipeline OCR et fuzzy matching.
\item
  Chapitre 4 --- Implémentation, Asynchronisme et Tests : interfaces,
  threading, notifications WebSocket, sécurité et validation.
\end{itemize}

\hypertarget{chapitre-1-contexte-du-projet}{%
\section{Chapitre 1 : Contexte du
Projet}\label{chapitre-1-contexte-du-projet}}

\hypertarget{introduction}{%
\subsection{1.1 Introduction}\label{introduction}}

Ce chapitre présente le cadre général dans lequel
s\textquotesingle inscrit notre stage de fin d\textquotesingle études.
Nous y décrivons l\textquotesingle organisme d\textquotesingle accueil,
le contexte institutionnel, la problématique identifiée et les objectifs
du projet.

\hypertarget{pruxe9sentation-de-lorganisme-daccueil}{%
\subsection{1.2 Présentation de l\textquotesingle organisme
d\textquotesingle accueil}\label{pruxe9sentation-de-lorganisme-daccueil}}

\hypertarget{lensa-buxe9ni-mellal}{%
\subsubsection{1.2.1 L\textquotesingle ENSA Béni
Mellal}\label{lensa-buxe9ni-mellal}}

L\textquotesingle École Nationale des Sciences Appliquées de Béni Mellal
(ENSA BM) est un établissement public d\textquotesingle enseignement
supérieur rattaché à l\textquotesingle Université Sultan Moulay Slimane
(USMS). Créée en 2003 par décret ministériel n° 2-03-200, elle fait
partie du réseau des douze ENSA marocaines et forme des ingénieurs
d\textquotesingle État dans les domaines des sciences appliquées.

\begin{longtable}[]{@{}
  >{\raggedright\arraybackslash}p{(\columnwidth - 2\tabcolsep) * \real{0.4972}}
  >{\raggedright\arraybackslash}p{(\columnwidth - 2\tabcolsep) * \real{0.5028}}@{}}
\toprule\noalign{}
\begin{minipage}[b]{\linewidth}\raggedright
\textbf{Rubrique}
\end{minipage} & \begin{minipage}[b]{\linewidth}\raggedright
\textbf{Information}
\end{minipage} \\
\midrule\noalign{}
\endhead
\bottomrule\noalign{}
\endlastfoot
Dénomination & École Nationale des Sciences Appliquées de Béni Mellal \\
Sigle & ENSA BM \\
Université & Université Sultan Moulay Slimane (USMS) \\
Statut & Établissement public d\textquotesingle enseignement
supérieur \\
Création & 2003 --- Décret n° 2-03-200 \\
Effectif étudiants & \textasciitilde2 500 (2024-2025) \\
Enseignants & \textasciitilde120 enseignants-chercheurs \\
\end{longtable}

Table 1: Fiche signalétique de l\textquotesingle ENSA BM

\hypertarget{missions-et-offre-de-formation}{%
\subsubsection{1.2.2 Missions et offre de
formation}\label{missions-et-offre-de-formation}}

L\textquotesingle ENSA BM a pour mission de former des ingénieurs
d\textquotesingle État polyvalents, de développer la recherche
scientifique et l\textquotesingle innovation technologique, et de
contribuer au développement socio-économique de la région Béni
Mellal-Khénifra.

L\textquotesingle école propose quatre filières
d\textquotesingle ingénieur accessibles aux niveaux Bac+2 et Bac+3 :

\begin{longtable}[]{@{}
  >{\raggedright\arraybackslash}p{(\columnwidth - 4\tabcolsep) * \real{0.3333}}
  >{\raggedright\arraybackslash}p{(\columnwidth - 4\tabcolsep) * \real{0.3333}}
  >{\raggedright\arraybackslash}p{(\columnwidth - 4\tabcolsep) * \real{0.3333}}@{}}
\toprule\noalign{}
\begin{minipage}[b]{\linewidth}\raggedright
\textbf{Code}
\end{minipage} & \begin{minipage}[b]{\linewidth}\raggedright
\textbf{Filière}
\end{minipage} & \begin{minipage}[b]{\linewidth}\raggedright
\textbf{Places}
\end{minipage} \\
\midrule\noalign{}
\endhead
\bottomrule\noalign{}
\endlastfoot
TDI & Télécommunications et Développement Informatique & 30 \\
IACS & Informatique Appliquée et Cybersécurité & 30 \\
IAA & Ingénierie Agro-Alimentaire & 30 \\
G2ER & Génie de l\textquotesingle Énergie et Énergies Renouvelables &
30 \\
\end{longtable}

Table 2: Filières d\textquotesingle ingénieur de l\textquotesingle ENSA
BM

\hypertarget{service-de-la-scolarituxe9}{%
\subsubsection{1.2.3 Service de la
Scolarité}\label{service-de-la-scolarituxe9}}

Le service de la Scolarité gère l\textquotesingle ensemble des processus
administratifs liés aux étudiants : inscriptions, examens, stages,
diplomation et présélection des candidats au cycle ingénieur.
C\textquotesingle est le principal bénéficiaire et commanditaire de
notre plateforme.

\hypertarget{contexte-et-probluxe9matique}{%
\subsection{1.3 Contexte et
problématique}\label{contexte-et-probluxe9matique}}

\hypertarget{le-processus-actuel-de-pruxe9suxe9lection}{%
\subsubsection{1.3.1 Le processus actuel de
présélection}\label{le-processus-actuel-de-pruxe9suxe9lection}}

Le processus actuel de présélection à l\textquotesingle ENSA BM se
déroule en plusieurs étapes manuelles : (1) réception des dossiers
physiques ou par email, (2) vérification manuelle de la conformité des
pièces, (3) saisie des notes dans un fichier Excel, (4) tri et
classement manuel, (5) convocation à l\textquotesingle épreuve écrite,
et (6) saisie des notes d\textquotesingle épreuve et classement final.
Ce processus mobilise plusieurs agents pendant plusieurs semaines.

\hypertarget{probluxe8mes-identifiuxe9s}{%
\subsubsection{1.3.2 Problèmes
identifiés}\label{probluxe8mes-identifiuxe9s}}

\begin{itemize}
\item
  Traitement lent et chronophage : plusieurs semaines pour quelques
  centaines de dossiers
\item
  Risque élevé d\textquotesingle erreurs de saisie lors de la
  transcription manuelle des notes
\item
  Absence totale de traçabilité des décisions prises
\item
  Manque de transparence : les candidats n\textquotesingle ont aucune
  visibilité sur l\textquotesingle état de leur dossier
\item
  Doublons non détectés : difficulté d\textquotesingle identifier les
  candidatures multiples
\item
  Classement approximatif : calculs manuels sujets à erreur
\end{itemize}

\hypertarget{le-systuxe8me-lmd-marocain}{%
\subsubsection{1.3.3 Le système LMD
marocain}\label{le-systuxe8me-lmd-marocain}}

Le système LMD marocain structure les études supérieures en trois
cycles. Chaque semestre est évalué sur 20 points. Les mentions sont
attribuées selon les seuils suivants :

\begin{longtable}[]{@{}
  >{\raggedright\arraybackslash}p{(\columnwidth - 2\tabcolsep) * \real{0.5001}}
  >{\raggedright\arraybackslash}p{(\columnwidth - 2\tabcolsep) * \real{0.4999}}@{}}
\toprule\noalign{}
\begin{minipage}[b]{\linewidth}\raggedright
\textbf{Moyenne}
\end{minipage} & \begin{minipage}[b]{\linewidth}\raggedright
\textbf{Mention}
\end{minipage} \\
\midrule\noalign{}
\endhead
\bottomrule\noalign{}
\endlastfoot
≥ 16.00 & Très Bien \\
≥ 14.00 & Bien \\
≥ 12.00 & Assez Bien \\
≥ 10.00 & Passable \\
\end{longtable}

Notre plateforme doit impérativement respecter ces seuils pour le calcul
automatique des mentions et du scoring académique.

\hypertarget{objectifs-du-projet}{%
\subsection{1.4 Objectifs du projet}\label{objectifs-du-projet}}

\hypertarget{objectif-principal}{%
\subsubsection{1.4.1 Objectif principal}\label{objectif-principal}}

Concevoir et développer une plateforme web complète permettant
d\textquotesingle automatiser et de fiabiliser le processus de
présélection des candidats aux filières d\textquotesingle ingénieur de
l\textquotesingle ENSA BM.

\hypertarget{objectifs-spuxe9cifiques}{%
\subsubsection{1.4.2 Objectifs
spécifiques}\label{objectifs-spuxe9cifiques}}

\begin{itemize}
\item
  Automatisation du rejet des dossiers non conformes via un moteur de
  règles configurable
\item
  Scoring académique pondéré par filière avec fuzzy matching des
  matières
\item
  Classement automatique des candidats intégrant dossier (40\%) et
  épreuve écrite (60\%)
\item
  Notifications en temps réel via WebSocket et email SMTP
\item
  Traçabilité et audit complet du processus de présélection
\item
  Import des notes d\textquotesingle épreuve écrite au format Excel
\end{itemize}

\hypertarget{puxe9rimuxe8tre-fonctionnel}{%
\subsection{1.5 Périmètre
fonctionnel}\label{puxe9rimuxe8tre-fonctionnel}}

\begin{longtable}[]{@{}
  >{\raggedright\arraybackslash}p{(\columnwidth - 2\tabcolsep) * \real{0.5001}}
  >{\raggedright\arraybackslash}p{(\columnwidth - 2\tabcolsep) * \real{0.4999}}@{}}
\toprule\noalign{}
\begin{minipage}[b]{\linewidth}\raggedright
\textbf{Fonctionnalité}
\end{minipage} & \begin{minipage}[b]{\linewidth}\raggedright
\textbf{Inclus}
\end{minipage} \\
\midrule\noalign{}
\endhead
\bottomrule\noalign{}
\endlastfoot
Inscription et authentification JWT & Oui \\
Dépôt de dossier multi-étapes & Oui \\
Moteur de règles de rejet automatique & Oui \\
Scoring pondéré par filière & Oui \\
Dashboard responsable avec KPIs & Oui \\
Import notes Excel & Oui \\
Notifications WebSocket + SMTP & Oui \\
Export résultats Excel & Oui \\
Application mobile native & Non (hors scope) \\
Intégration Massar & Non (hors scope) \\
\end{longtable}

Table 3: Périmètre fonctionnel du projet

\hypertarget{muxe9thodologie-de-travail}{%
\subsection{1.6 Méthodologie de
travail}\label{muxe9thodologie-de-travail}}

Nous avons adopté une approche agile itérative, organisée en sprints
hebdomadaires sur 4 semaines. Chaque sprint produit un incrément
fonctionnel testé et validé.

\begin{longtable}[]{@{}
  >{\raggedright\arraybackslash}p{(\columnwidth - 4\tabcolsep) * \real{0.3333}}
  >{\raggedright\arraybackslash}p{(\columnwidth - 4\tabcolsep) * \real{0.3333}}
  >{\raggedright\arraybackslash}p{(\columnwidth - 4\tabcolsep) * \real{0.3333}}@{}}
\toprule\noalign{}
\begin{minipage}[b]{\linewidth}\raggedright
\textbf{Semaine}
\end{minipage} & \begin{minipage}[b]{\linewidth}\raggedright
\textbf{Phase}
\end{minipage} & \begin{minipage}[b]{\linewidth}\raggedright
\textbf{Livrables}
\end{minipage} \\
\midrule\noalign{}
\endhead
\bottomrule\noalign{}
\endlastfoot
S1 & Fondations Backend & Modèles Django, API REST, Auth JWT \\
S2 & Intelligence métier & Moteur de règles, Scoring, OCR \\
S3 & Frontend React & Pages candidat, Dashboard, Formulaires \\
S4 & Tests et finalisation & Tests, Corrections, Documentation \\
\end{longtable}

Table 4: Planning prévisionnel du projet

\hypertarget{conclusion}{%
\subsection{1.7 Conclusion}\label{conclusion}}

Ce chapitre a permis de situer le cadre institutionnel de notre stage et
de définir clairement la problématique et les objectifs du projet. Le
chapitre suivant présentera l\textquotesingle architecture technique et
le modèle de données de la plateforme.

\hypertarget{chapitre-2-architecture-et-base-de-donnuxe9es}{%
\section{Chapitre 2 : Architecture et Base de
données}\label{chapitre-2-architecture-et-base-de-donnuxe9es}}

\hypertarget{introduction-1}{%
\subsection{2.1 Introduction}\label{introduction-1}}

Ce chapitre détaille l\textquotesingle architecture technique de la
plateforme, le choix des technologies, l\textquotesingle organisation
modulaire du code et le dictionnaire complet des modèles Django.

\includegraphics[width=6.18264in,height=6.79722in]{media/image3.png}

\hypertarget{architecture-globale}{%
\subsection{2.2 Architecture globale}\label{architecture-globale}}

\hypertarget{vue-densemble}{%
\subsubsection{2.2.1 Vue
d\textquotesingle ensemble}\label{vue-densemble}}

L\textquotesingle application repose sur une architecture moderne
Full-Stack Client-Serveur à trois niveaux :

\begin{itemize}
\item
  Frontend : React 18 compilé via Vite, avec Tailwind CSS pour le rendu
  responsive. L\textquotesingle architecture s\textquotesingle appuie
  sur des composants réutilisables, React Router v6 et un state
  management optimisé via Context API.
\item
  Backend : Django 4.2 LTS et Django REST Framework (DRF), exposant une
  API REST sécurisée par JWT (SimpleJWT). Le backend est organisé en 5
  applications Django indépendantes.
\item
  Base de données : PostgreSQL 15 géré par l\textquotesingle ORM Django,
  assurant la persistance des candidatures, configurations et
  historiques.
\item
  Temps réel : Django Channels 4.0 avec Redis comme channel layer pour
  les WebSockets (notifications en temps réel avec barre de
  progression).

  \includegraphics[width=2.81319in,height=4.18125in]{media/image5.png}
\end{itemize}

\hypertarget{architecture-rest-stateless}{%
\subsubsection{2.2.2 Architecture REST
stateless}\label{architecture-rest-stateless}}

Le frontend communique avec le backend exclusivement via des endpoints
REST sécurisés par JWT. Les réponses sont au format JSON. Le pattern est
strictement stateless côté serveur, à l\textquotesingle exception des
connexions WebSocket persistantes pour les notifications.

\hypertarget{stack-technologique}{%
\subsection{2.3 Stack technologique}\label{stack-technologique}}

\begin{longtable}[]{@{}
  >{\raggedright\arraybackslash}p{(\columnwidth - 4\tabcolsep) * \real{0.3333}}
  >{\raggedright\arraybackslash}p{(\columnwidth - 4\tabcolsep) * \real{0.3333}}
  >{\raggedright\arraybackslash}p{(\columnwidth - 4\tabcolsep) * \real{0.3333}}@{}}
\toprule\noalign{}
\begin{minipage}[b]{\linewidth}\raggedright
\textbf{Technologie}
\end{minipage} & \begin{minipage}[b]{\linewidth}\raggedright
\textbf{Version}
\end{minipage} & \begin{minipage}[b]{\linewidth}\raggedright
\textbf{Rôle}
\end{minipage} \\
\midrule\noalign{}
\endhead
\bottomrule\noalign{}
\endlastfoot
Django & 4.2 LTS & Framework backend principal \\
Django REST Framework & 3.17 & Construction de l\textquotesingle API
REST \\
Django Channels & 4.0 & WebSockets temps réel \\
SimpleJWT & 5.5 & Authentification par tokens JWT \\
React & 18 & Framework frontend SPA \\
Tailwind CSS & 3.x & Framework CSS utilitaire \\
PostgreSQL & 15 & Base de données relationnelle \\
Redis & 7.x & Channel Layer WebSocket \\
pdfplumber + pytesseract & --- & Extraction OCR \\
fuzzywuzzy & 0.18 & Correspondance floue des matières \\
openpyxl & --- & Import/Export Excel \\
\end{longtable}

Table 5: Stack technologique du projet

\hypertarget{organisation-modulaire-du-backend}{%
\subsection{2.4 Organisation modulaire du
backend}\label{organisation-modulaire-du-backend}}

Le backend est structuré en 5 applications Django indépendantes, chacune
gérant un domaine métier spécifique :

\begin{longtable}[]{@{}
  >{\raggedright\arraybackslash}p{(\columnwidth - 4\tabcolsep) * \real{0.3333}}
  >{\raggedright\arraybackslash}p{(\columnwidth - 4\tabcolsep) * \real{0.3333}}
  >{\raggedright\arraybackslash}p{(\columnwidth - 4\tabcolsep) * \real{0.3333}}@{}}
\toprule\noalign{}
\begin{minipage}[b]{\linewidth}\raggedright
\textbf{Application}
\end{minipage} & \begin{minipage}[b]{\linewidth}\raggedright
\textbf{Responsabilité}
\end{minipage} & \begin{minipage}[b]{\linewidth}\raggedright
\textbf{Fichiers clés}
\end{minipage} \\
\midrule\noalign{}
\endhead
\bottomrule\noalign{}
\endlastfoot
users & Authentification, comptes, profils candidats & models.py,
serializers.py, views.py \\
candidatures & Dossiers, documents, notes, services métier & models.py,
views.py, services/ \\
administration & Filières, diplômes, épreuves, historique & models.py,
views.py \\
scoring & Règles de rejet, configuration scoring & models.py,
views.py \\
notifications & Notifications email et WebSocket & models.py,
consumers.py, views.py \\
\end{longtable}

Table 6: Applications Django du backend

Le répertoire candidatures/services/ contient la logique métier critique
:

\begin{longtable}[]{@{}
  >{\raggedright\arraybackslash}p{(\columnwidth - 4\tabcolsep) * \real{0.3333}}
  >{\raggedright\arraybackslash}p{(\columnwidth - 4\tabcolsep) * \real{0.3333}}
  >{\raggedright\arraybackslash}p{(\columnwidth - 4\tabcolsep) * \real{0.3333}}@{}}
\toprule\noalign{}
\begin{minipage}[b]{\linewidth}\raggedright
\textbf{Service}
\end{minipage} & \begin{minipage}[b]{\linewidth}\raggedright
\textbf{Fichier}
\end{minipage} & \begin{minipage}[b]{\linewidth}\raggedright
\textbf{Rôle}
\end{minipage} \\
\midrule\noalign{}
\endhead
\bottomrule\noalign{}
\endlastfoot
Extraction OCR & extraction.py & Extraction des informations depuis les
documents scannés \\
Fuzzy Matching & fuzzy\_matching.py & Correspondance floue des noms de
matières \\
Règles de rejet & regles.py & Application des 7 types de règles de rejet
automatique \\
Scoring & scoring.py & Calcul du score pondéré par filière \\
Vérification & verification\_document.py & Vérification
d\textquotesingle authenticité des documents \\
Invitation PDF & invitation\_pdf.py & Génération des convocations au
format PDF \\
\end{longtable}

Table 7: Services métier du module candidatures

\includegraphics[width=2.35417in,height=3.96875in]{media/image6.png}

\hypertarget{dictionnaire-des-moduxe8les-django}{%
\subsection{2.5 Dictionnaire des modèles
Django}\label{dictionnaire-des-moduxe8les-django}}

Le modèle relationnel est structuré autour de 12 entités principales
réparties dans les applications Django. Nous détaillons ci-dessous
chaque modèle.

\hypertarget{application-users}{%
\subsubsection{2.5.1 Application Users}\label{application-users}}

Le modèle User étend AbstractBaseUser de Django pour personnaliser
l\textquotesingle authentification. L\textquotesingle identification se
fait par email et CIN (Carte d\textquotesingle Identité Nationale).
Trois rôles sont définis : CANDIDAT, RESPONSABLE et ADMIN, implémentant
un contrôle d\textquotesingle accès RBAC.

\begin{longtable}[]{@{}
  >{\raggedright\arraybackslash}p{(\columnwidth - 4\tabcolsep) * \real{0.3333}}
  >{\raggedright\arraybackslash}p{(\columnwidth - 4\tabcolsep) * \real{0.3333}}
  >{\raggedright\arraybackslash}p{(\columnwidth - 4\tabcolsep) * \real{0.3333}}@{}}
\toprule\noalign{}
\begin{minipage}[b]{\linewidth}\raggedright
\textbf{Attribut}
\end{minipage} & \begin{minipage}[b]{\linewidth}\raggedright
\textbf{Type}
\end{minipage} & \begin{minipage}[b]{\linewidth}\raggedright
\textbf{Contrainte}
\end{minipage} \\
\midrule\noalign{}
\endhead
\bottomrule\noalign{}
\endlastfoot
id & UUID & Clé primaire auto-générée \\
email & VARCHAR & UNIQUE, NOT NULL \\
cin & VARCHAR & UNIQUE, NOT NULL \\
role & ENUM & CANDIDAT / RESPONSABLE / ADMIN \\
is\_active & BOOLEAN & Défaut : True \\
password & VARCHAR & Hashé via PBKDF2 \\
\end{longtable}

Le modèle Candidat est relié en OneToOne à User et stocke le profil
détaillé :

\begin{longtable}[]{@{}
  >{\raggedright\arraybackslash}p{(\columnwidth - 4\tabcolsep) * \real{0.3333}}
  >{\raggedright\arraybackslash}p{(\columnwidth - 4\tabcolsep) * \real{0.3333}}
  >{\raggedright\arraybackslash}p{(\columnwidth - 4\tabcolsep) * \real{0.3333}}@{}}
\toprule\noalign{}
\begin{minipage}[b]{\linewidth}\raggedright
\textbf{Attribut}
\end{minipage} & \begin{minipage}[b]{\linewidth}\raggedright
\textbf{Type}
\end{minipage} & \begin{minipage}[b]{\linewidth}\raggedright
\textbf{Contrainte}
\end{minipage} \\
\midrule\noalign{}
\endhead
\bottomrule\noalign{}
\endlastfoot
id & UUID & Clé primaire \\
user\_id & FK → User & OneToOne, CASCADE \\
nom & VARCHAR & NOT NULL \\
prenom & VARCHAR & NOT NULL \\
telephone & VARCHAR & Optionnel \\
date\_naissance & DATE & NOT NULL \\
\end{longtable}

\hypertarget{application-candidatures}{%
\subsubsection{2.5.2 Application
Candidatures}\label{application-candidatures}}

Le modèle Dossier est le modèle central de
l\textquotesingle application. Il lie un Candidat à une Filière et
contient le statut, les scores calculés et les motifs de rejet
éventuels.

\begin{longtable}[]{@{}
  >{\raggedright\arraybackslash}p{(\columnwidth - 4\tabcolsep) * \real{0.3333}}
  >{\raggedright\arraybackslash}p{(\columnwidth - 4\tabcolsep) * \real{0.3333}}
  >{\raggedright\arraybackslash}p{(\columnwidth - 4\tabcolsep) * \real{0.3333}}@{}}
\toprule\noalign{}
\begin{minipage}[b]{\linewidth}\raggedright
\textbf{Attribut}
\end{minipage} & \begin{minipage}[b]{\linewidth}\raggedright
\textbf{Type}
\end{minipage} & \begin{minipage}[b]{\linewidth}\raggedright
\textbf{Contrainte}
\end{minipage} \\
\midrule\noalign{}
\endhead
\bottomrule\noalign{}
\endlastfoot
id & UUID & Clé primaire \\
candidat\_id & FK → Candidat & CASCADE \\
filiere\_id & FK → Filiere & CASCADE \\
statut & ENUM & 11 valeurs possibles \\
score & DECIMAL & Score dossier calculé \\
score\_final & DECIMAL & 40\% dossier + 60\% épreuve \\
rang\_final & INT & Classement par filière \\
mention & ENUM & Auto-calculé \\
moyenne\_generale & DECIMAL & Moyenne des semestres \\
\end{longtable}

Le modèle Document gère les pièces justificatives uploadées par le
candidat :

\begin{longtable}[]{@{}
  >{\raggedright\arraybackslash}p{(\columnwidth - 4\tabcolsep) * \real{0.3333}}
  >{\raggedright\arraybackslash}p{(\columnwidth - 4\tabcolsep) * \real{0.3333}}
  >{\raggedright\arraybackslash}p{(\columnwidth - 4\tabcolsep) * \real{0.3333}}@{}}
\toprule\noalign{}
\begin{minipage}[b]{\linewidth}\raggedright
\textbf{Attribut}
\end{minipage} & \begin{minipage}[b]{\linewidth}\raggedright
\textbf{Type}
\end{minipage} & \begin{minipage}[b]{\linewidth}\raggedright
\textbf{Contrainte}
\end{minipage} \\
\midrule\noalign{}
\endhead
\bottomrule\noalign{}
\endlastfoot
id & UUID & Clé primaire \\
dossier\_id & FK → Dossier & CASCADE \\
type\_doc & ENUM & DIPLOME / RELEVE / CIN / PHOTO \\
fichier & FILE & Stocké dans media/dossiers/\{uuid\}/ \\
qualite\_ocr & DECIMAL & Score de confiance OCR \\
\end{longtable}

\includegraphics[width=7.84653in,height=3.66875in]{media/image7.png}Le
modèle NoteSemestre remplace l\textquotesingle ancienne approche par
NoteMatiere. Il intègre la moyenne par semestre, la session et la
mention calculée dynamiquement :

\begin{longtable}[]{@{}
  >{\raggedright\arraybackslash}p{(\columnwidth - 4\tabcolsep) * \real{0.3333}}
  >{\raggedright\arraybackslash}p{(\columnwidth - 4\tabcolsep) * \real{0.3333}}
  >{\raggedright\arraybackslash}p{(\columnwidth - 4\tabcolsep) * \real{0.3333}}@{}}
\toprule\noalign{}
\begin{minipage}[b]{\linewidth}\raggedright
\textbf{Attribut}
\end{minipage} & \begin{minipage}[b]{\linewidth}\raggedright
\textbf{Type}
\end{minipage} & \begin{minipage}[b]{\linewidth}\raggedright
\textbf{Contrainte}
\end{minipage} \\
\midrule\noalign{}
\endhead
\bottomrule\noalign{}
\endlastfoot
id & UUID & Clé primaire \\
dossier\_id & FK → Dossier & CASCADE \\
semestre & ENUM & S1 à S6 \\
moyenne & DECIMAL & Plage stricte {[}10.00, 20.00{]} \\
session & ENUM & NORMALE / RATTRAPAGE \\
mention & ENUM & Auto-calculé (editable=False) \\
\end{longtable}

\hypertarget{application-administration-et-scoring}{%
\subsubsection{2.5.3 Application Administration et
Scoring}\label{application-administration-et-scoring}}

Le modèle Filiere définit les filières du cycle ingénieur avec leurs
paramètres :

\begin{longtable}[]{@{}
  >{\raggedright\arraybackslash}p{(\columnwidth - 4\tabcolsep) * \real{0.3333}}
  >{\raggedright\arraybackslash}p{(\columnwidth - 4\tabcolsep) * \real{0.3333}}
  >{\raggedright\arraybackslash}p{(\columnwidth - 4\tabcolsep) * \real{0.3333}}@{}}
\toprule\noalign{}
\begin{minipage}[b]{\linewidth}\raggedright
\textbf{Attribut}
\end{minipage} & \begin{minipage}[b]{\linewidth}\raggedright
\textbf{Type}
\end{minipage} & \begin{minipage}[b]{\linewidth}\raggedright
\textbf{Contrainte}
\end{minipage} \\
\midrule\noalign{}
\endhead
\bottomrule\noalign{}
\endlastfoot
id & UUID & Clé primaire \\
nom & VARCHAR & NOT NULL \\
code & VARCHAR & UNIQUE (TDI, IACS, IAA, G2ER) \\
niveau & ENUM & BAC2 / BAC3 \\
places\_disponibles & INT & Défaut : 30 \\
\end{longtable}

Les modèles EpreuveEcrite et NoteEcrite gèrent la logistique du concours
écrit :

\begin{longtable}[]{@{}
  >{\raggedright\arraybackslash}p{(\columnwidth - 4\tabcolsep) * \real{0.3333}}
  >{\raggedright\arraybackslash}p{(\columnwidth - 4\tabcolsep) * \real{0.3333}}
  >{\raggedright\arraybackslash}p{(\columnwidth - 4\tabcolsep) * \real{0.3333}}@{}}
\toprule\noalign{}
\begin{minipage}[b]{\linewidth}\raggedright
\textbf{Modèle}
\end{minipage} & \begin{minipage}[b]{\linewidth}\raggedright
\textbf{Attributs clés}
\end{minipage} & \begin{minipage}[b]{\linewidth}\raggedright
\textbf{Rôle}
\end{minipage} \\
\midrule\noalign{}
\endhead
\bottomrule\noalign{}
\endlastfoot
EpreuveEcrite & nom, date\_epreuve, seuil\_admission, statut & Gestion
de l\textquotesingle épreuve par filière \\
NoteEcrite & note, resultat, rang\_final & Note individuelle et
classement \\
ConfigScoring & matiere, poids, bonus\_mention & Pondération scoring par
filière \\
RegleRejet & type\_regle, parametre, is\_active, message & Règle de
rejet configurable \\
HistoriqueAction & action, commentaire, timestamp & Traçabilité des
décisions \\
\end{longtable}

Table 8: Modèles d\textquotesingle administration et scoring

\hypertarget{relations-et-cardinalituxe9s}{%
\subsection{2.6 Relations et
cardinalités}\label{relations-et-cardinalituxe9s}}

\begin{longtable}[]{@{}
  >{\raggedright\arraybackslash}p{(\columnwidth - 4\tabcolsep) * \real{0.3333}}
  >{\raggedright\arraybackslash}p{(\columnwidth - 4\tabcolsep) * \real{0.3333}}
  >{\raggedright\arraybackslash}p{(\columnwidth - 4\tabcolsep) * \real{0.3333}}@{}}
\toprule\noalign{}
\begin{minipage}[b]{\linewidth}\raggedright
\textbf{Relation}
\end{minipage} & \begin{minipage}[b]{\linewidth}\raggedright
\textbf{Type}
\end{minipage} & \begin{minipage}[b]{\linewidth}\raggedright
\textbf{Cardinalité}
\end{minipage} \\
\midrule\noalign{}
\endhead
\bottomrule\noalign{}
\endlastfoot
User → Candidat & OneToOne & 1..1 \\
Candidat → Dossier & OneToMany & 1..* \\
Dossier → Document & OneToMany & 1..4 \\
Dossier → NoteSemestre & OneToMany & 1..6 \\
Filiere → Dossier & OneToMany & 1..* \\
Filiere → RegleRejet & OneToMany & 1..* \\
Filiere → EpreuveEcrite & OneToMany & 1..* \\
EpreuveEcrite → NoteEcrite & OneToMany & 1..* \\
Dossier → HistoriqueAction & OneToMany & 1..* \\
User → Notification & OneToMany & 1..* \\
\end{longtable}

Table 9: Relations et cardinalités du modèle de données

\hypertarget{machine-duxe9tats-du-dossier}{%
\subsection{2.7 Machine d\textquotesingle états du
dossier}\label{machine-duxe9tats-du-dossier}}

Le dossier suit une machine d\textquotesingle états stricte implémentée
dans la méthode changer\_statut() du modèle. Chaque transition crée
automatiquement un enregistrement HistoriqueAction pour la traçabilité :

\begin{longtable}[]{@{}
  >{\raggedright\arraybackslash}p{(\columnwidth - 2\tabcolsep) * \real{0.5001}}
  >{\raggedright\arraybackslash}p{(\columnwidth - 2\tabcolsep) * \real{0.4999}}@{}}
\toprule\noalign{}
\begin{minipage}[b]{\linewidth}\raggedright
\textbf{État}
\end{minipage} & \begin{minipage}[b]{\linewidth}\raggedright
\textbf{Transitions possibles}
\end{minipage} \\
\midrule\noalign{}
\endhead
\bottomrule\noalign{}
\endlastfoot
BROUILLON & → EN\_TRAITEMENT (soumission) \\
EN\_TRAITEMENT & → REJETE\_AUTO \textbar{} EN\_ATTENTE \textbar{}
INCOMPLET \textbar{} SUSPECT \\
INCOMPLET & → EN\_TRAITEMENT (complétion) \\
SUSPECT & → EN\_ATTENTE \textbar{} REJETE\_FINAL (décision manuelle) \\
EN\_ATTENTE & → PRESELECTIONNE \textbar{} REJETE\_FINAL \\
PRESELECTIONNE & → ADMIS\_FINAL \textbar{} RECALE\_FINAL \textbar{}
ABSENT\_ECRIT \\
\end{longtable}

Table 10: Tableau 2.6 : Machine d\textquotesingle états du dossier de
candidature

\includegraphics[width=4.29236in,height=6.62847in]{media/image9.png}

\hypertarget{conclusion-1}{%
\subsection{2.8 Conclusion}\label{conclusion-1}}

Ce chapitre a détaillé l\textquotesingle architecture Full-Stack de la
plateforme et le modèle de données relationnel complet.
L\textquotesingle organisation modulaire en 5 applications Django assure
une séparation claire des responsabilités. Le chapitre suivant
présentera la logique métier et les algorithmes de scoring et
d\textquotesingle OCR.

\hypertarget{chapitre-3-logique-muxe9tier-et-algorithmes}{%
\section{Chapitre 3 : Logique Métier et
Algorithmes}\label{chapitre-3-logique-muxe9tier-et-algorithmes}}

\hypertarget{introduction-2}{%
\subsection{3.1 Introduction}\label{introduction-2}}

Ce chapitre expose en détail les algorithmes et la logique métier qui
constituent le cœur intelligent de la plateforme. Nous y décrivons le
moteur de règles de rejet, l\textquotesingle algorithme de scoring
pondéré par filière, le pipeline OCR avec fuzzy matching, et les
formules mathématiques appliquées.

\hypertarget{moteur-de-ruxe8gles-de-rejet-automatique}{%
\subsection{3.2 Moteur de règles de rejet
automatique}\label{moteur-de-ruxe8gles-de-rejet-automatique}}

\hypertarget{principe-de-fonctionnement}{%
\subsubsection{3.2.1 Principe de
fonctionnement}\label{principe-de-fonctionnement}}

Le moteur de règles applique séquentiellement un ensemble de
vérifications configurables par filière. Dès qu\textquotesingle une
règle est violée, le dossier est automatiquement rejeté avec le motif
correspondant (pattern court-circuit). Chaque règle peut être activée ou
désactivée individuellement par l\textquotesingle administrateur.

\hypertarget{types-de-ruxe8gles-impluxe9mentuxe9s}{%
\subsubsection{3.2.2 Types de règles
implémentés}\label{types-de-ruxe8gles-impluxe9mentuxe9s}}

\begin{longtable}[]{@{}
  >{\raggedright\arraybackslash}p{(\columnwidth - 4\tabcolsep) * \real{0.3333}}
  >{\raggedright\arraybackslash}p{(\columnwidth - 4\tabcolsep) * \real{0.3333}}
  >{\raggedright\arraybackslash}p{(\columnwidth - 4\tabcolsep) * \real{0.3333}}@{}}
\toprule\noalign{}
\begin{minipage}[b]{\linewidth}\raggedright
\textbf{Code Règle}
\end{minipage} & \begin{minipage}[b]{\linewidth}\raggedright
\textbf{Description}
\end{minipage} & \begin{minipage}[b]{\linewidth}\raggedright
\textbf{Paramètre}
\end{minipage} \\
\midrule\noalign{}
\endhead
\bottomrule\noalign{}
\endlastfoot
DIPLOME\_INVALIDE & Diplôme non accepté pour la filière cible & Liste
des diplômes autorisés \\
MOYENNE\_INSUFFISANTE & Moyenne générale en dessous du seuil minimal &
Seuil minimal (ex: 10.00) \\
NOTE\_ELIMINATOIRE & Note éliminatoire dans une matière clé & Seuil et
matière \\
DOUBLON\_CIN & CIN déjà utilisé dans un autre dossier actif &
Automatique \\
DOCUMENT\_MANQUANT & Document obligatoire absent du dossier & Liste des
types requis \\
DATE\_INCOHERENTE & Incohérence dans les dates du diplôme &
Automatique \\
ETABLISSEMENT\_INVALIDE & Établissement d\textquotesingle origine non
reconnu & Liste blanche \\
\end{longtable}

Table 11: Les 7 types de règles de rejet automatique

\includegraphics[width=7.12153in,height=1.84861in]{media/image11.png}

\hypertarget{journalisation-des-rejets}{%
\subsubsection{3.2.3 Journalisation des
rejets}\label{journalisation-des-rejets}}

Chaque rejet automatique génère un enregistrement HistoriqueAction
contenant le motif précis de rejet, la règle déclenchée et
l\textquotesingle horodatage. Le statut du dossier passe à REJETE\_AUTO
et le candidat est notifié automatiquement par email.

\hypertarget{algorithme-de-scoring-ponduxe9ruxe9-par-filiuxe8re}{%
\subsection{3.3 Algorithme de scoring pondéré par
filière}\label{algorithme-de-scoring-ponduxe9ruxe9-par-filiuxe8re}}

\hypertarget{principe-guxe9nuxe9ral}{%
\subsubsection{3.3.1 Principe général}\label{principe-guxe9nuxe9ral}}

Le scoring repose sur une pondération configurable par filière. Chaque
filière définit des coefficients par catégorie de matières. Les moyennes
semestrielles du candidat sont regroupées par catégorie via le fuzzy
matching, puis pondérées selon les coefficients de la filière cible.

\hypertarget{coefficients-des-quatre-filiuxe8res}{%
\subsubsection{3.3.2 Coefficients des quatre
filières}\label{coefficients-des-quatre-filiuxe8res}}

Le moteur de scoring applique les pondérations suivantes sur les
catégories de matières extraites, en fonction de la filière choisie par
le candidat :

\begin{longtable}[]{@{}
  >{\raggedright\arraybackslash}p{(\columnwidth - 12\tabcolsep) * \real{0.1429}}
  >{\raggedright\arraybackslash}p{(\columnwidth - 12\tabcolsep) * \real{0.1428}}
  >{\raggedright\arraybackslash}p{(\columnwidth - 12\tabcolsep) * \real{0.1430}}
  >{\raggedright\arraybackslash}p{(\columnwidth - 12\tabcolsep) * \real{0.1428}}
  >{\raggedright\arraybackslash}p{(\columnwidth - 12\tabcolsep) * \real{0.1431}}
  >{\raggedright\arraybackslash}p{(\columnwidth - 12\tabcolsep) * \real{0.1427}}
  >{\raggedright\arraybackslash}p{(\columnwidth - 12\tabcolsep) * \real{0.1429}}@{}}
\toprule\noalign{}
\begin{minipage}[b]{\linewidth}\raggedright
\textbf{Filière}
\end{minipage} & \begin{minipage}[b]{\linewidth}\raggedright
\textbf{INFO}
\end{minipage} & \begin{minipage}[b]{\linewidth}\raggedright
\textbf{MATH}
\end{minipage} & \begin{minipage}[b]{\linewidth}\raggedright
\textbf{ELEC\_AUTO}
\end{minipage} & \begin{minipage}[b]{\linewidth}\raggedright
\textbf{CHIMIE\_BIO}
\end{minipage} & \begin{minipage}[b]{\linewidth}\raggedright
\textbf{RESEAUX\_BD}
\end{minipage} & \begin{minipage}[b]{\linewidth}\raggedright
\textbf{LANGUES}
\end{minipage} \\
\midrule\noalign{}
\endhead
\bottomrule\noalign{}
\endlastfoot
TDI & 35\% & 30\% & 25\% & --- & --- & 10\% \\
IACS & 40\% & 35\% & --- & --- & 15\% & 10\% \\
IAA & 10\% & 20\% & --- & 60\% & --- & 10\% \\
G2ER & 15\% & 30\% & 25\% & 30\% & --- & --- \\
\end{longtable}

Table 12: Matrice des pondérations par filière

\hypertarget{formule-de-calcul-du-score-dossier}{%
\subsubsection{3.3.3 Formule de calcul du score
dossier}\label{formule-de-calcul-du-score-dossier}}

Pour chaque filière F, le score dossier S\_d est calculé comme suit :

\textbf{S\_d(F) = Σ {[} w\_i(F) × M\_i {]}}

Où :

\begin{itemize}
\item
  w\_i(F) : poids de la catégorie de matière i pour la filière F
\item
  M\_i : moyenne du candidat dans la catégorie i (obtenue par fuzzy
  matching)
\item
  Si une catégorie indispensable n\textquotesingle est pas détectée, le
  moteur applique la moyenne générale comme fallback et marque le
  dossier comme SUSPECT

  \includegraphics[width=2.97361in,height=6.12569in]{media/image13.png}
\end{itemize}

\hypertarget{bonus-de-mention-et-coefficient-dipluxf4me}{%
\subsubsection{3.3.4 Bonus de mention et coefficient
diplôme}\label{bonus-de-mention-et-coefficient-dipluxf4me}}

Le score brut est enrichi par deux mécanismes de bonification :

\begin{longtable}[]{@{}
  >{\raggedright\arraybackslash}p{(\columnwidth - 2\tabcolsep) * \real{0.5001}}
  >{\raggedright\arraybackslash}p{(\columnwidth - 2\tabcolsep) * \real{0.4999}}@{}}
\toprule\noalign{}
\begin{minipage}[b]{\linewidth}\raggedright
\textbf{Mention}
\end{minipage} & \begin{minipage}[b]{\linewidth}\raggedright
\textbf{Bonus}
\end{minipage} \\
\midrule\noalign{}
\endhead
\bottomrule\noalign{}
\endlastfoot
Très Bien (≥ 16) & +2.0 points \\
Bien (≥ 14) & +1.5 points \\
Assez Bien (≥ 12) & +1.0 point \\
Passable (≥ 10) & +0.0 point \\
\end{longtable}

Un coefficient multiplicateur est appliqué selon le type de diplôme :

\begin{longtable}[]{@{}
  >{\raggedright\arraybackslash}p{(\columnwidth - 2\tabcolsep) * \real{0.5001}}
  >{\raggedright\arraybackslash}p{(\columnwidth - 2\tabcolsep) * \real{0.4999}}@{}}
\toprule\noalign{}
\begin{minipage}[b]{\linewidth}\raggedright
\textbf{Type de diplôme}
\end{minipage} & \begin{minipage}[b]{\linewidth}\raggedright
\textbf{Coefficient}
\end{minipage} \\
\midrule\noalign{}
\endhead
\bottomrule\noalign{}
\endlastfoot
DUT / BTS (spécialité correspondante) & ×1.00 \\
Licence (spécialité correspondante) & ×1.00 \\
Autre diplôme reconnu & ×0.80 \\
\end{longtable}

\hypertarget{formule-du-score-final}{%
\subsubsection{3.3.5 Formule du score
final}\label{formule-du-score-final}}

Le score final intègre la note de l\textquotesingle épreuve écrite selon
la formule :

\textbf{Score\_final = (Score\_dossier × 0.40) +
(Note\_épreuve\_normalisée × 0.60)}

En cas d\textquotesingle égalité de score final, le départage se fait
par : (1) la note d\textquotesingle épreuve écrite, (2) la moyenne
générale du dossier, (3) la date de soumission.

\hypertarget{pipeline-ocr-et-fuzzy-matching}{%
\subsection{3.4 Pipeline OCR et Fuzzy
Matching}\label{pipeline-ocr-et-fuzzy-matching}}

\hypertarget{extraction-ocr}{%
\subsubsection{3.4.1 Extraction OCR}\label{extraction-ocr}}

Le pipeline de traitement automatique des relevés de notes
s\textquotesingle appuie sur les bibliothèques pdfplumber (pour les PDF
natifs) et pytesseract (pour les documents scannés). Le texte brut est
extrait puis soumis à un processus de nettoyage.

\includegraphics[width=2.71319in,height=4.03056in]{media/image15.png}

\hypertarget{nettoyage-du-bruit-ocr}{%
\subsubsection{3.4.2 Nettoyage du bruit
OCR}\label{nettoyage-du-bruit-ocr}}

Une fonction spécialisée filtre les erreurs fréquentes
d\textquotesingle OCR sur les chiffres et les caractères :

\begin{longtable}[]{@{}
  >{\raggedright\arraybackslash}p{(\columnwidth - 2\tabcolsep) * \real{0.5001}}
  >{\raggedright\arraybackslash}p{(\columnwidth - 2\tabcolsep) * \real{0.4999}}@{}}
\toprule\noalign{}
\begin{minipage}[b]{\linewidth}\raggedright
\textbf{Caractère OCR}
\end{minipage} & \begin{minipage}[b]{\linewidth}\raggedright
\textbf{Correction}
\end{minipage} \\
\midrule\noalign{}
\endhead
\bottomrule\noalign{}
\endlastfoot
O (lettre) & 0 (chiffre) \\
l (lettre L min) & 1 (chiffre) \\
S (lettre) & 5 (chiffre) \\
I (lettre I maj) & 1 (chiffre) \\
\end{longtable}

Table 13: Corrections automatiques du bruit OCR

Les textes sont également normalisés : suppression des accents, passage
en minuscules, suppression des caractères spéciaux.

\hypertarget{fuzzy-matching-des-matiuxe8res}{%
\subsubsection{3.4.3 Fuzzy Matching des
matières}\label{fuzzy-matching-des-matiuxe8res}}

L\textquotesingle algorithme utilise la bibliothèque fuzzywuzzy pour
comparer les noms de matières extraits par OCR avec un dictionnaire de
catégories standardisées. Le processus :

\begin{itemize}
\item
  Chaque nom de matière extrait est comparé à toutes les entrées du
  dictionnaire
\item
  La fonction fuzz.token\_sort\_ratio calcule un score de similarité
  entre 0 et 100
\item
  Un seuil strict de 80\% est appliqué : seules les correspondances
  au-dessus sont retenues
\item
  Les catégories standardisées sont : MATH, INFO, CHIMIE\_BIO,
  ELEC\_AUTO, RESEAUX\_BD, LANGUES
\end{itemize}

Exemple : "Mathématiques pour l\textquotesingle ingénieur" → MATH
(similarité 92\%)

Exemple : "Algorithmique et structures de données" → INFO (similarité
85\%)

\hypertarget{scoring-avec-fallback}{%
\subsubsection{3.4.4 Scoring avec
fallback}\label{scoring-avec-fallback}}

Les notes extraites sont regroupées par catégorie standardisée. Si une
catégorie indispensable à la filière n\textquotesingle est pas détectée
par le fuzzy matching, le moteur applique la moyenne générale du
candidat comme valeur de remplacement (fallback) et marque
automatiquement le dossier comme SUSPECT pour vérification manuelle par
le responsable.

\hypertarget{calcul-automatique-des-mentions}{%
\subsection{3.5 Calcul automatique des
mentions}\label{calcul-automatique-des-mentions}}

La mention est calculée automatiquement dans la méthode save() du modèle
NoteSemestre. La validation stricte impose que les moyennes soient dans
la plage {[}10.00, 20.00{]}. Toute valeur hors plage déclenche une
ValidationError Django.

Le champ mention est déclaré avec editable=False pour empêcher toute
modification manuelle. Il est recalculé à chaque sauvegarde selon les
seuils LMD définis précédemment dans la section 1.3.3.

\hypertarget{conclusion-2}{%
\subsection{3.6 Conclusion}\label{conclusion-2}}

Ce chapitre a présenté les algorithmes au cœur de la plateforme : le
moteur de règles de rejet à 7 types, le scoring pondéré par filière avec
ses formules mathématiques, et le pipeline OCR avec fuzzy matching. Ces
mécanismes garantissent un traitement automatisé, fiable et transparent
des candidatures. Le chapitre suivant détaillera
l\textquotesingle implémentation, la gestion de
l\textquotesingle asynchronisme et les tests.

\hypertarget{chapitre-4-impluxe9mentation-asynchronisme-et-tests}{%
\section{Chapitre 4 : Implémentation, Asynchronisme et
Tests}\label{chapitre-4-impluxe9mentation-asynchronisme-et-tests}}

\hypertarget{introduction-3}{%
\subsection{4.1 Introduction}\label{introduction-3}}

Ce chapitre présente les aspects concrets de
l\textquotesingle implémentation : les interfaces développées, la
gestion de l\textquotesingle asynchronisme pour les opérations
coûteuses, le système de notifications en temps réel, la sécurité et la
stratégie de tests.

\hypertarget{interfaces-duxe9veloppuxe9es}{%
\subsection{4.2 Interfaces
développées}\label{interfaces-duxe9veloppuxe9es}}

L\textquotesingle application comprend 11 interfaces principales,
couvrant les trois rôles d\textquotesingle utilisateurs (Candidat,
Responsable, Administrateur) :

\begin{longtable}[]{@{}
  >{\raggedright\arraybackslash}p{(\columnwidth - 4\tabcolsep) * \real{0.3333}}
  >{\raggedright\arraybackslash}p{(\columnwidth - 4\tabcolsep) * \real{0.3333}}
  >{\raggedright\arraybackslash}p{(\columnwidth - 4\tabcolsep) * \real{0.3333}}@{}}
\toprule\noalign{}
\begin{minipage}[b]{\linewidth}\raggedright
\textbf{Interface}
\end{minipage} & \begin{minipage}[b]{\linewidth}\raggedright
\textbf{Rôle}
\end{minipage} & \begin{minipage}[b]{\linewidth}\raggedright
\textbf{Description}
\end{minipage} \\
\midrule\noalign{}
\endhead
\bottomrule\noalign{}
\endlastfoot
Page d\textquotesingle accueil & Public & Landing page avec statistiques
ENSA BM et filières \\
Authentification & Public & Inscription (email + CIN) et connexion
JWT \\
Formulaire candidature & Candidat & 5 étapes avec sauvegarde
brouillon \\
Suivi de dossier & Candidat & Statut temps réel avec badges colorés \\
Résultats épreuve & Candidat & Note, rang et résultat (Admis/Recalé) \\
Dashboard responsable & Responsable & KPIs, graphiques Recharts,
actions \\
Gestion des dossiers & Responsable & Liste paginée, filtres, tri par
score \\
Détail dossier & Responsable & Vue complète avec historique et
alertes \\
Configuration admin & Admin & Filières, règles, pondérations \\
Import Excel & Responsable & Wizard 3 étapes avec mapping colonnes \\
Notifications & Tous & Cloche, panneau déroulant, historique \\
\end{longtable}

Table 14: Interfaces développées par rôle

\includegraphics[width=6.33472in,height=2.97917in]{media/image17.png}

\hypertarget{section-1}{%
\subsection{}\label{section-1}}

\includegraphics[width=6.33472in,height=2.97917in]{media/image18.png}

\hypertarget{section-2}{%
\subsection{}\label{section-2}}

\includegraphics[width=6.33472in,height=2.97917in]{media/image19.png}

\hypertarget{section-3}{%
\subsection{\texorpdfstring{\protect\includegraphics[width=6.33472in,height=2.97917in]{media/image20.png}}{}}\label{section-3}}

\hypertarget{section-4}{%
\subsection{\texorpdfstring{\protect\includegraphics[width=6.33472in,height=2.97917in]{media/image21.png}}{}}\label{section-4}}

\hypertarget{section-5}{%
\subsection{\texorpdfstring{\protect\includegraphics[width=6.33472in,height=3.53333in]{media/image22.png}}{}}\label{section-5}}

\hypertarget{section-6}{%
\subsection{\texorpdfstring{\protect\includegraphics[width=6.33472in,height=3.50278in]{media/image23.png}}{}}\label{section-6}}

\hypertarget{gestion-de-lasynchronisme}{%
\subsection{4.3 Gestion de
l\textquotesingle asynchronisme}\label{gestion-de-lasynchronisme}}

\hypertarget{probluxe9matique}{%
\subsubsection{4.3.1 Problématique}\label{probluxe9matique}}

Plusieurs opérations de la plateforme sont coûteuses en temps
d\textquotesingle exécution : le traitement OCR des documents, le calcul
des scores, l\textquotesingle envoi d\textquotesingle emails en masse.
Un traitement synchrone bloquerait l\textquotesingle interface
utilisateur pendant plusieurs secondes.

\hypertarget{solution-threading.thread}{%
\subsubsection{4.3.2 Solution :
threading.Thread}\label{solution-threading.thread}}

Pour garantir une excellente expérience utilisateur,
l\textquotesingle application utilise l\textquotesingle asynchronisme
via threading.Thread de Python. Les processus lourds sont exécutés en
arrière-plan dans des threads démons. Le flux est le suivant :

\begin{itemize}
\item
  Le candidat soumet son dossier → réponse HTTP 200 immédiate
\item
  Un thread démon est lancé pour le traitement OCR et le scoring
\item
  Le thread ferme correctement sa connexion à la base de données après
  exécution
\item
  Le candidat est redirigé instantanément vers son dashboard
\item
  Le statut du dossier est mis à jour en arrière-plan et notifié via
  WebSocket
\end{itemize}

\hypertarget{gestion-des-connexions-base-de-donnuxe9es}{%
\subsubsection{4.3.3 Gestion des connexions base de
données}\label{gestion-des-connexions-base-de-donnuxe9es}}

Le système veille à clore correctement les connexions à la base de
données dans les threads via django.db.connections.close\_all() pour
éviter toute fuite de connexion ou conflit d\textquotesingle accès
concurrent.

\hypertarget{systuxe8me-de-notifications-en-temps-ruxe9el}{%
\subsection{4.4 Système de notifications en temps
réel}\label{systuxe8me-de-notifications-en-temps-ruxe9el}}

\hypertarget{architecture-websocket}{%
\subsubsection{4.4.1 Architecture
WebSocket}\label{architecture-websocket}}

Les notifications combinent deux canaux complémentaires :

\begin{itemize}
\item
  SMTP (Gmail App Password) : envoi d\textquotesingle emails avec
  gestion des erreurs et statut persisté
\item
  Django Channels + Redis : notifications WebSocket temps réel avec
  barre de progression pour les envois en masse
\end{itemize}

\hypertarget{envoi-en-masse-avec-progression}{%
\subsubsection{4.4.2 Envoi en masse avec
progression}\label{envoi-en-masse-avec-progression}}

Lors de l\textquotesingle envoi de notifications en masse, le backend
envoie les emails par batch SMTP et notifie le frontend de la
progression en temps réel via WebSocket. Le responsable voit une barre
de progression se remplir au fur et à mesure des envois.

\hypertarget{systuxe8me-de-logging}{%
\subsection{4.5 Système de logging}\label{systuxe8me-de-logging}}

L\textquotesingle application intègre un mécanisme de traçabilité
sécurisé par fichiers. La configuration Django dirige les flux vers un
fichier central (logs/ensa\_presel.log). Les logs sont structurés par
domaine :

\begin{longtable}[]{@{}
  >{\raggedright\arraybackslash}p{(\columnwidth - 4\tabcolsep) * \real{0.3333}}
  >{\raggedright\arraybackslash}p{(\columnwidth - 4\tabcolsep) * \real{0.3333}}
  >{\raggedright\arraybackslash}p{(\columnwidth - 4\tabcolsep) * \real{0.3333}}@{}}
\toprule\noalign{}
\begin{minipage}[b]{\linewidth}\raggedright
\textbf{Logger}
\end{minipage} & \begin{minipage}[b]{\linewidth}\raggedright
\textbf{Niveau}
\end{minipage} & \begin{minipage}[b]{\linewidth}\raggedright
\textbf{Usage}
\end{minipage} \\
\midrule\noalign{}
\endhead
\bottomrule\noalign{}
\endlastfoot
django & WARNING & Erreurs framework \\
candidatures & INFO & Opérations métier \\
notifications & INFO & Envois email/WebSocket \\
channels & DEBUG & Connexions WebSocket \\
\end{longtable}

Table 15: Configuration des loggers

\hypertarget{suxe9curituxe9-de-lapplication}{%
\subsection{4.6 Sécurité de
l\textquotesingle application}\label{suxe9curituxe9-de-lapplication}}

\begin{longtable}[]{@{}
  >{\raggedright\arraybackslash}p{(\columnwidth - 2\tabcolsep) * \real{0.5001}}
  >{\raggedright\arraybackslash}p{(\columnwidth - 2\tabcolsep) * \real{0.4999}}@{}}
\toprule\noalign{}
\begin{minipage}[b]{\linewidth}\raggedright
\textbf{Mesure}
\end{minipage} & \begin{minipage}[b]{\linewidth}\raggedright
\textbf{Implémentation}
\end{minipage} \\
\midrule\noalign{}
\endhead
\bottomrule\noalign{}
\endlastfoot
Authentification & JWT (access 15min + refresh 24h) \\
Autorisation & RBAC à 3 niveaux via @permission\_classes \\
CSRF & Protection Django middleware activée \\
XSS & Échappement automatique Django + React \\
Injection SQL & ORM Django avec requêtes paramétrées \\
Logs sécurisés & Fichier avec rotation, pas de console \\
Confidentialité & Scores OCR masqués côté candidat \\
\end{longtable}

Table 16: Mesures de sécurité implémentées

\hypertarget{stratuxe9gie-de-tests}{%
\subsection{4.7 Stratégie de tests}\label{stratuxe9gie-de-tests}}

\hypertarget{niveaux-de-tests}{%
\subsubsection{4.7.1 Niveaux de tests}\label{niveaux-de-tests}}

\begin{longtable}[]{@{}
  >{\raggedright\arraybackslash}p{(\columnwidth - 4\tabcolsep) * \real{0.3333}}
  >{\raggedright\arraybackslash}p{(\columnwidth - 4\tabcolsep) * \real{0.3333}}
  >{\raggedright\arraybackslash}p{(\columnwidth - 4\tabcolsep) * \real{0.3333}}@{}}
\toprule\noalign{}
\begin{minipage}[b]{\linewidth}\raggedright
\textbf{Niveau}
\end{minipage} & \begin{minipage}[b]{\linewidth}\raggedright
\textbf{Objectif}
\end{minipage} & \begin{minipage}[b]{\linewidth}\raggedright
\textbf{Outils}
\end{minipage} \\
\midrule\noalign{}
\endhead
\bottomrule\noalign{}
\endlastfoot
Tests unitaires & Valider chaque composant isolément & Django
TestCase \\
Tests d\textquotesingle intégration & Valider les flux complets &
APITestCase, APIClient \\
Tests fonctionnels & Valider les scénarios utilisateur & Scénarios
manuels \\
Tests d\textquotesingle interface & Valider le responsive et
l\textquotesingle UX & Navigateur, DevTools \\
\end{longtable}

\hypertarget{tests-du-moteur-de-ruxe8gles}{%
\subsubsection{4.7.2 Tests du moteur de
règles}\label{tests-du-moteur-de-ruxe8gles}}

\begin{longtable}[]{@{}
  >{\raggedright\arraybackslash}p{(\columnwidth - 6\tabcolsep) * \real{0.2500}}
  >{\raggedright\arraybackslash}p{(\columnwidth - 6\tabcolsep) * \real{0.2500}}
  >{\raggedright\arraybackslash}p{(\columnwidth - 6\tabcolsep) * \real{0.2500}}
  >{\raggedright\arraybackslash}p{(\columnwidth - 6\tabcolsep) * \real{0.2499}}@{}}
\toprule\noalign{}
\begin{minipage}[b]{\linewidth}\raggedright
\textbf{Test}
\end{minipage} & \begin{minipage}[b]{\linewidth}\raggedright
\textbf{Entrée}
\end{minipage} & \begin{minipage}[b]{\linewidth}\raggedright
\textbf{Résultat attendu}
\end{minipage} & \begin{minipage}[b]{\linewidth}\raggedright
\textbf{Statut}
\end{minipage} \\
\midrule\noalign{}
\endhead
\bottomrule\noalign{}
\endlastfoot
Diplôme invalide & DUT non accepté pour TDI & REJETE\_AUTO & ✓ Passé \\
Moyenne insuffisante & Moyenne 9.5 (seuil 10) & REJETE\_AUTO & ✓
Passé \\
Doublon CIN & CIN déjà existant & REJETE\_AUTO & ✓ Passé \\
Dossier conforme & Tous critères OK & EN\_ATTENTE & ✓ Passé \\
\end{longtable}

\hypertarget{tests-de-lalgorithme-de-scoring}{%
\subsubsection{4.7.3 Tests de l\textquotesingle algorithme de
scoring}\label{tests-de-lalgorithme-de-scoring}}

\begin{longtable}[]{@{}
  >{\raggedright\arraybackslash}p{(\columnwidth - 6\tabcolsep) * \real{0.2500}}
  >{\raggedright\arraybackslash}p{(\columnwidth - 6\tabcolsep) * \real{0.2500}}
  >{\raggedright\arraybackslash}p{(\columnwidth - 6\tabcolsep) * \real{0.2500}}
  >{\raggedright\arraybackslash}p{(\columnwidth - 6\tabcolsep) * \real{0.2499}}@{}}
\toprule\noalign{}
\begin{minipage}[b]{\linewidth}\raggedright
\textbf{Test}
\end{minipage} & \begin{minipage}[b]{\linewidth}\raggedright
\textbf{Données}
\end{minipage} & \begin{minipage}[b]{\linewidth}\raggedright
\textbf{Score attendu}
\end{minipage} & \begin{minipage}[b]{\linewidth}\raggedright
\textbf{Statut}
\end{minipage} \\
\midrule\noalign{}
\endhead
\bottomrule\noalign{}
\endlastfoot
Scoring TDI & Moyennes S1-S6, DUT Info & 15.42 & ✓ Passé \\
Scoring IACS & Moyennes S1-S6, Licence Info & 14.87 & ✓ Passé \\
Bonus mention TB & Moyenne ≥16 & +2.0 bonus & ✓ Passé \\
\end{longtable}

\hypertarget{ruxe9capitulatif-des-tests}{%
\subsubsection{4.7.4 Récapitulatif des
tests}\label{ruxe9capitulatif-des-tests}}

\begin{longtable}[]{@{}
  >{\raggedright\arraybackslash}p{(\columnwidth - 6\tabcolsep) * \real{0.2500}}
  >{\raggedright\arraybackslash}p{(\columnwidth - 6\tabcolsep) * \real{0.2500}}
  >{\raggedright\arraybackslash}p{(\columnwidth - 6\tabcolsep) * \real{0.2500}}
  >{\raggedright\arraybackslash}p{(\columnwidth - 6\tabcolsep) * \real{0.2499}}@{}}
\toprule\noalign{}
\begin{minipage}[b]{\linewidth}\raggedright
\textbf{Type de test}
\end{minipage} & \begin{minipage}[b]{\linewidth}\raggedright
\textbf{Nombre}
\end{minipage} & \begin{minipage}[b]{\linewidth}\raggedright
\textbf{Passés}
\end{minipage} & \begin{minipage}[b]{\linewidth}\raggedright
\textbf{Échoués}
\end{minipage} \\
\midrule\noalign{}
\endhead
\bottomrule\noalign{}
\endlastfoot
Unitaires & 32 & 32 & 0 \\
Intégration & 12 & 12 & 0 \\
Fonctionnels & 9 & 9 & 0 \\
Interface & 8 & 8 & 0 \\
Total & 61 & 61 & 0 \\
\end{longtable}

\emph{Tableau 4.4 : Récapitulatif des résultats de tests}

\hypertarget{couverture-de-code}{%
\subsubsection{4.7.5 Couverture de code}\label{couverture-de-code}}

\begin{longtable}[]{@{}
  >{\raggedright\arraybackslash}p{(\columnwidth - 2\tabcolsep) * \real{0.5001}}
  >{\raggedright\arraybackslash}p{(\columnwidth - 2\tabcolsep) * \real{0.4999}}@{}}
\toprule\noalign{}
\begin{minipage}[b]{\linewidth}\raggedright
\textbf{Module}
\end{minipage} & \begin{minipage}[b]{\linewidth}\raggedright
\textbf{Couverture}
\end{minipage} \\
\midrule\noalign{}
\endhead
\bottomrule\noalign{}
\endlastfoot
users & 89\% \\
candidatures & 92\% \\
administration & 85\% \\
scoring & 94\% \\
notifications & 78\% \\
Global & 87\% \\
\end{longtable}

\hypertarget{difficultuxe9s-rencontruxe9es-et-solutions}{%
\subsection{4.8 Difficultés rencontrées et
solutions}\label{difficultuxe9s-rencontruxe9es-et-solutions}}

\begin{longtable}[]{@{}
  >{\raggedright\arraybackslash}p{(\columnwidth - 2\tabcolsep) * \real{0.5001}}
  >{\raggedright\arraybackslash}p{(\columnwidth - 2\tabcolsep) * \real{0.4999}}@{}}
\toprule\noalign{}
\begin{minipage}[b]{\linewidth}\raggedright
\textbf{Difficulté}
\end{minipage} & \begin{minipage}[b]{\linewidth}\raggedright
\textbf{Solution apportée}
\end{minipage} \\
\midrule\noalign{}
\endhead
\bottomrule\noalign{}
\endlastfoot
WebSocket + JWT expiré & Middleware WebSocket avec refresh automatique
du token \\
SMTP Gmail bloqué & Configuration App Password avec gestion des
erreurs \\
Imports circulaires Django & Import local dans les méthodes au lieu
d\textquotesingle imports globaux \\
OCR faible qualité & Fallback sur moyenne générale + flag SUSPECT \\
Connexions DB dans threads & Appel systématique à
connections.close\_all() \\
\end{longtable}

Table 17: Difficultés rencontrées et solutions

\hypertarget{conclusion-3}{%
\subsection{4.9 Conclusion}\label{conclusion-3}}

Ce chapitre a présenté les aspects concrets de
l\textquotesingle implémentation : les 11 interfaces développées, la
gestion de l\textquotesingle asynchronisme via threading, le système de
notifications WebSocket temps réel, la sécurité robuste et la stratégie
de tests complète avec un taux de réussite de 100\% et une couverture de
code de 87\%.

\hypertarget{conclusion-guxe9nuxe9rale}{%
\section{Conclusion générale}\label{conclusion-guxe9nuxe9rale}}

Le stage effectué au sein de l\textquotesingle ENSA Béni Mellal nous a
permis de concevoir et développer une plateforme web complète de
présélection des candidats aux filières d\textquotesingle ingénieur. Ce
projet a couvert l\textquotesingle ensemble du cycle de développement
logiciel : analyse des besoins, conception architecturale,
implémentation Full Stack et tests.

La plateforme développée répond aux objectifs fixés : automatisation du
traitement des dossiers via un moteur de règles configurable à 7 types,
scoring académique pondéré par filière avec fuzzy matching des matières,
gestion des épreuves écrites avec import Excel, notifications en temps
réel via WebSocket, et traçabilité complète des décisions.

Sur le plan technique, l\textquotesingle architecture Django REST
Framework + React 18 s\textquotesingle est révélée robuste et
performante. Les 5 applications Django assurent une séparation claire
des responsabilités, tandis que le frontend React offre une expérience
utilisateur fluide et responsive. La stratégie de tests a permis
d\textquotesingle atteindre un taux de réussite de 100\% sur 61 tests
avec une couverture de code de 87\%.

Ce stage a été une expérience enrichissante qui m\textquotesingle a
permis de consolider mes compétences en développement Full Stack, en
architecture logicielle et en gestion de projet agile.

\hypertarget{perspectives-damuxe9lioration}{%
\subsection{Perspectives
d\textquotesingle amélioration}\label{perspectives-damuxe9lioration}}

\begin{itemize}
\item
  Intégration avec la plateforme nationale Massar via convention
  ministérielle pour la vérification automatique des CNE
\item
  Développement d\textquotesingle une application mobile (Flutter/React
  Native) pour le suivi
\item
  Tableau de bord analytique avancé avec indicateurs de performance
  détaillés
\item
  Intelligence artificielle pour la détection de fraude documentaire
  avancée
\item
  Déploiement en production avec conteneurisation Docker et CI/CD
\end{itemize}

\hypertarget{bibliographie}{%
\section{Bibliographie}\label{bibliographie}}

{[}1{]} Django Software Foundation, « Django Documentation v4.2 LTS »,
docs.djangoproject.com, 2024.

{[}2{]} Meta Platforms, « React Documentation v18 », react.dev, 2024.

{[}3{]} T. Christie, « Django REST Framework Documentation »,
django-rest-framework.org, 2024.

{[}4{]} J. Davie, « Simple JWT Documentation »,
django-rest-framework-simplejwt.readthedocs.io, 2024.

{[}5{]} Django Channels, « Channels Documentation v4.0 »,
channels.readthedocs.io, 2024.

{[}6{]} Tailwind Labs, « Tailwind CSS Documentation v3 »,
tailwindcss.com, 2024.

{[}7{]} PostgreSQL Global Development Group, « PostgreSQL 15
Documentation », postgresql.org, 2024.

{[}8{]} Redis Ltd, « Redis Documentation v7 », redis.io, 2024.

{[}9{]} pdfplumber, « Documentation », github.com/jsvine/pdfplumber,
2024.

{[}10{]} SeatGeek, « fuzzywuzzy --- Fuzzy String Matching in Python »,
github.com/seatgeek/fuzzywuzzy, 2024.

{[}11{]} C. Larman, « UML 2 et les Design Patterns », Pearson Education,
3e édition, 2005.

{[}12{]} Ministère de l\textquotesingle Enseignement Supérieur, «
Système LMD au Maroc », enssup.gov.ma, 2023.

{[}13{]} ENSA Béni Mellal, « Site officiel », ensabm.usms.ac.ma, 2024.

\end{document}
